% Part: first-order-logic
% Chapter: completeness
% Section: compactness

\documentclass[../../../include/open-logic-section]{subfiles}

\begin{document}

\iftag{FOL}
      {\olfileid[pt]{fol}{com}{com}}
      {\olfileid[pt]{pl}{com}{com}}

\olsection{O Teorema da Compacidade}

Uma consequência importante do teorema da completude é o
teorema da compacidade.  O teorema da compacidade afirma que, se cada
subconjunto \emph{finito} de um conjunto de !!{sentence}s é satisfatível, o
conjunto inteiro é satisfatível --- mesmo que o próprio conjunto seja infinito. Isso
está longe de ser óbvio. Não há nada que pareça descartar, ao menos à primeira
vista, a possibilidade de existirem conjuntos infinitos de
!!{sentence}s que sejam contraditórios, mas em que a contradição só
surge, por assim dizer, do número infinito.  O teorema da compacidade
diz que tal cenário pode ser descartado: não existem
conjuntos infinitos insatisfatíveis de !!{sentence}s tais que cada subconjunto finito
deles seja satisfatível. Como o teorema da completude, ele tem uma versão
relacionada à consequência lógica: se um conjunto infinito de !!{sentence}s implica
algo, então já um subconjunto finito o implica.

\begin{defn}
  Um conjunto $\Gamma$ de !!{formula}s é \emph{finitamente satisfatível} se e somente se todo $\Gamma_0 \subseteq \Gamma$ finito é satisfatível.
\end{defn}

\begin{thm}[Teorema da Compacidade]
\ollabel{thm:compactness}
O seguinte vale para quaisquer sentenças $\Gamma$ e $!A$:
\begin{enumerate}
  \item $\Gamma \Entails !A$ se e somente se existe um $\Gamma_0
    \subseteq \Gamma$ finito tal que $\Gamma_0 \Entails !A$.
  \item $\Gamma$ é satisfatível se e somente se é finitamente
    satisfatível.
\end{enumerate}
\end{thm}

\begin{proof}
Provamos (2).  Se $\Gamma$ é satisfatível, então existe
\iftag{FOL}{!!a{structure}~$\Struct{M}$}{!!a{valuation}~$\pAssign{v}$}
tal que $\iftag{FOL}{\Sat{M}}{\pSat{v}}{!A}$ para todo $!A \in
\Gamma$.  É claro que essa $\iftag{FOL}{\Struct{M}}{\pAssign{v}}$ também
satisfaz todo subconjunto finito de~$\Gamma$, de modo que $\Gamma$ é finitamente
satisfatível.

Agora suponha que $\Gamma$ seja finitamente satisfatível.  Então todo subconjunto
finito~$\Gamma_0 \subseteq \Gamma$ é satisfatível.  Pela correção
(\iftag{FOL}{%
  \tagrefs{prfAX/{fol:axd:sou:cor:consistency-soundness},
    prfSC/{fol:seq:sou:cor:consistency-soundness},
    prfND/{fol:ntd:sou:cor:consistency-soundness},
    prfTab/{fol:tab:sou:cor:consistency-soundness}}}{%
  \tagrefs{prfAX/{pl:axd:sou:cor:consistency-soundness},
    prfSC/{pl:seq:sou:cor:consistency-soundness},
    prfND/{pl:ntd:sou:cor:consistency-soundness},
    prfTab/{pl:tab:sou:cor:consistency-soundness}}}),
todo subconjunto finito é consistente.  Então o próprio $\Gamma$ deve ser
consistente por
\iftag{FOL}{%
  \tagrefs{prfAX/{fol:axd:ptn:prop:proves-compact},
    prfSC/{fol:seq:ptn:prop:proves-compact},
    prfND/{fol:ntd:ptn:prop:proves-compact},
    prfTab/{fol:tab:ptn:prop:proves-compact}}}{%
  \tagrefs{prfAX/{pl:axd:ptn:prop:proves-compact},
    prfSC/{pl:seq:ptn:prop:proves-compact},
    prfND/{pl:ntd:ptn:prop:proves-compact},
    prfTab/{pl:tab:ptn:prop:proves-compact}}}.
Pela completude (\olref[cth]{thm:completeness}), como $\Gamma$~é
consistente, ele é satisfatível.
\end{proof}

\tagprob{FOL}
\begin{prob}
Prove (1) de \olref[fol][com][com]{thm:compactness}.
\end{prob}
\tagendprob

\tagprob{notFOL}
\begin{prob}
Prove (1) de \olref[pl][com][com]{thm:compactness}.
\end{prob}
\tagendprob

\iftag{FOL}{%
\begin{ex}
Em todo modelo~$\Struct{M}$ de uma teoria~$\Gamma$, cada termo~$t$
obviamente designa !!a{element} de~$\Domain{M}$. Podemos garantir que
também seja verdade que todo !!{element} de~$\Domain{M}$ seja designado
por algum termo? Em outras palavras, existem teorias~$\Gamma$ cujos
modelos são todos cobertos?  O teorema da compacidade mostra que esse
não é o caso se $\Gamma$ tem modelos infinitos. Eis como ver
isso: Seja $\Struct{M}$ um modelo infinito de~$\Gamma$, e seja $c$
!!a{constant} que não está na linguagem de~$\Gamma$. Seja $\Delta$ o
conjunto de todas as sentenças $\eq/[c][t]$ para $t$ um termo na
linguagem~$\Lang{L}$ de~$\Gamma$, isto é,
  \[
  \Delta = \Setabs{\eq/[c][t]}{t \in \Trm[L]}.
  \]
Um subconjunto finito de $\Gamma \cup \Delta$ pode ser escrito como $\Gamma'
\cup \Delta'$, com $\Gamma' \subseteq \Gamma$ e $\Delta' \subseteq
\Delta$. Como $\Delta'$ é finito, ele pode conter apenas finitos
termos. Seja $a \in \Domain{M}$ !!a{element} de $\Domain{M}$ que não é
designado por nenhum deles, e seja $\Struct{M'}$ a !!{structure}
que é exatamente como $\Struct{M}$, mas com $\Assign{c}{M'} = a$ também. Como
$a \neq \Value{t}{M}$ para todo~$t$ que ocorre em~$\Delta'$,
$\Sat{M'}{\Delta'}$. Como $\Sat{M}{\Gamma}$, $\Gamma' \subseteq
\Gamma$, e $c$ não ocorre em~$\Gamma$, também
$\Sat{M'}{\Gamma'}$. Juntando, $\Sat{M'}{\Gamma' \cup \Delta'}$ para
todo subconjunto finito $\Gamma' \cup \Delta'$ de $\Gamma \cup \Delta$. Assim,
todo subconjunto finito de $\Gamma \cup \Delta$ é satisfatível. Por
compacidade, o próprio $\Gamma \cup \Delta$ é satisfatível. Logo, existem
modelos~$\Sat{M}{\Gamma \cup \Delta}$. Todo tal $\Struct{M}$ é um
modelo de~$\Gamma$, mas não é coberto, já que $\Value{c}{M} \neq
\Value{t}{M}$ para todos os termos~$t$ de~$\Lang{L}$.
\end{ex}

\begin{ex}
Considere !!a{language} $\Lang{L}$ contendo o !!{predicate}~$<$,
!!{constant}s $\Obj{0}$, $\Obj{1}$, e !!{function}s $+$, $\times$, e
$-$. Seja $\Gamma$ o conjunto de todas as !!{sentence}s nesta
!!{language} verdadeiras na !!{structure}~$\Struct{Q}$ com domínio~$\Rat$ e as
interpretações óbvias.  $\Gamma$~é o conjunto de todas as !!{sentence}s
de~$\Lang{L}$ verdadeiras a respeito dos números racionais. É claro que, em~$\Rat$
(e mesmo em~$\Real$), não há números~$r$ que sejam maiores que~$0$
mas menores que $1/k$ para todo $k \in \PosInt$.  Tal número, se
existisse, seria um \emph{infinitesimal:} não nulo, mas infinitamente
pequeno.  O teorema da compacidade pode ser usado para mostrar que há
modelos de~$\Gamma$ nos quais infinitesimais existem. Não temos
!!a{function} para a divisão em nossa linguagem (a divisão por zero é
indefinida, e !!{function}s têm de ser interpretadas por funções totais).
Contudo, ainda podemos expressar que $r < 1/k$, já que isso ocorre se e somente se
$r \cdot k < 1$. Agora seja $c$ !!a{constant} nova e seja $\Delta$
\[
\{\Obj{0} < c\}
\cup \Setabs{ c \times \num k < \Obj{1} }{k \in \PosInt}
\]
(onde $\num{k} = (\Obj{1} + (\Obj{1} + \dots + (\Obj{1} +
\Obj{1})\dots))$ com $k$~$\Obj{1}$'s). Para qualquer subconjunto
finito~$\Delta_0$ de~$\Delta$ existe um~$K$ tal que todas as
!!{sentence}s $c \times \num{k} < \Obj{1}$ em~$\Delta_0$ têm $k < K$.
Se expandirmos $\Struct{Q}$ para~$\Struct{Q'}$ com $\Assign{c}{Q'} = 1/K$
temos que $\Sat{Q'}{\Gamma_0 \cup \Delta_0}$ para qualquer
$\Gamma_0 \subseteq \Gamma$ finito, e assim $\Gamma \cup \Delta$ é finitamente
satisfatível (Exercício: prove isso em detalhe). Por compacidade, $\Gamma
\cup \Delta$ é satisfatível. Qualquer modelo~$\Struct{S}$ de $\Gamma \cup
\Delta$ contém um infinitesimal, a saber~$\Assign{c}{S}$.
\end{ex}
}{}

\tagprob{FOL}
\begin{prob}
No modelo padrão da aritmética~$\Struct{N}$, não há
!!{element}~$k \in \Domain{N}$ que satisfaça toda fórmula~$\num{n}
< x$ (onde $\num{n}$ é $\Obj{0}^{\prime\dots\prime}$ com $n$
$\prime$'s).  Use o teorema da compacidade para mostrar que o conjunto de
!!{sentence}s na linguagem da aritmética que são verdadeiras no modelo padrão
da aritmética $\Struct{N}$ também são verdadeiras em
!!a{structure}~$\Struct{N'}$ que contém !!a{element} que
\emph{de fato} satisfaz toda fórmula $\num{n} < x$.
\end{prob}
\tagendprob

\iftag{FOL}{
\begin{ex}
Sabemos que a lógica de primeira ordem com !!{identity} pode expressar que o
tamanho do !!{domain} deve ter algum tamanho mínimo: A
sentença~$!A_{\ge n}$ (que diz ``existem pelo menos $n$ objetos
distintos'') é verdadeira apenas em estruturas onde $\Domain{M}$ tem pelo
menos~$n$ objetos. Assim, se tomarmos
  \[
  \Delta = \Setabs{!A_{\ge n}}{n \ge 1}
  \]
então qualquer modelo de $\Delta$ deve ser infinito. Desse modo, podemos garantir
que uma teoria só tenha modelos infinitos adicionando~$\Delta$ a ela: os
modelos de $\Gamma \cup \Delta$ são todos os modelos infinitos
de~$\Gamma$, e apenas eles.

Portanto, a lógica de primeira ordem pode expressar a infinitude. O teorema da compacidade
mostra que ela não pode expressar a finitude, no entanto. Pois suponha que algum conjunto
de sentenças $\Lambda$ fosse satisfeito em todas as !!{structure}s finitas, e
apenas nelas. Então $\Delta \cup \Lambda$ é finitamente
satisfatível. Por quê? Suponha que $\Delta' \cup \Lambda' \subseteq \Delta \cup
\Lambda$ seja finito, com $\Delta' \subseteq \Delta$ e $\Lambda'
\subseteq \Lambda$. Seja $n$ o maior número tal que $!A_{\ge n}
\in \Delta'$. $\Lambda$, sendo satisfeito em todas as estruturas finitas, tem
um modelo~$\Struct{M}$ com um número finito de !!{element}s, mas $\ge n$ deles. Mas
então $\Sat{M}{\Delta' \cup \Lambda'}$.  Por compacidade, $\Delta \cup
\Lambda$ tem um modelo infinito, contradizendo a suposição de que
$\Lambda$ é satisfeito apenas em !!{structure}s finitas.
\end{ex}
}{}

\end{document}
